% Source: Wald GR, physical PDF pp. 393--402
%         (printed pp. 380--389).
% Coverage: Section 14.1, Quantum Gravity; equations (14.1.1)--(14.1.10).
% Figures/tables/footnotes: no figures or tables; footnotes 1--5.
% Status: source-reviewed and compile-clean.
% Modernization: all script symbols use \mathcal. The source perturbation
%                \gamma_{ab} is written \delta g_{ab} rather than h_{ab} here
%                to avoid collision with the spatial metrics (h_1)_{ab},
%                (h_2)_{ab} later in the same discussion, as NOTATION.md directs.
% Uncertainties: none.

\subsection{Quantum Gravity}\label{sec:14.1}

It is generally believed that the correct, fundamental description of all physical
fields is given by the general framework of quantum field theory. In quantum field
theory, states of a system are described by vectors in a Hilbert space \(\mathcal H\), and the
physical field is described by an operator (i.e., a linear map) on \(\mathcal H\) defined at each
spacetime point. However, unlike ordinary, nonrelativistic Schrödinger quantum
mechanics where well defined---albeit probabilistic---predictions always can be
made once one is given the Hamiltonian of the system, serious difficulties arise when
one attempts to formulate quantum field theories. Many of these difficulties can be
traced to the fact that, even for a free field, the expression obtained for the field
operator does not make mathematical sense as an operator defined at each spacetime
point but must be interpreted as a distribution on spacetime (see section~14.2 below).
This corresponds to the physical fact that the field cannot be measured at a single
point; only averages of the field over spacetime regions are physically well defined.
In the case of a free field this causes no serious problems and a well defined quantum
field theory can be constructed. However, for the more interesting case of a theory
with interactions (i.e., for a field or fields satisfying nonlinear equations) one is led
unavoidably to consider the products of field operators at the same spacetime point.
Such quantities have no natural mathematical meaning, and consequently, with the
exception of some simple models in lower spacetime dimensions, there are at present
no known examples of quantum field theories of physically reasonable interacting
fields which are mathematically well defined.

However, although one does not know how to formulate an exact theory of an
interacting quantum field, one can formally treat the interactions as perturbations of
the well defined free field theory. One thereby may obtain a formal expansion for
physical quantities as power series in the ``coupling constant'' (i.e., the coefficient of
the nonlinear interaction term in the equation). However, in general the formal
expression for the individual terms in the power series will yield divergent results
when one attempts to evaluate them, as should be expected from the fact that the
exact theory upon which the perturbation series is based is ill defined. Nevertheless,
one can introduce cutoffs into the divergent expressions in order to obtain finite
answers. This, of course, leaves the theory in an unsatisfactory state, since the
predicted values of all physical quantities depend upon the chosen values of the
cutoff parameters. However, it may happen that in the limit of large cutoff parameters,
the dependence of all physical quantities on the cutoff parameters may be
identical to their dependence on the so-called bare parameters (such as masses,
charges, and the coupling constants) which originally were present in the theory. If
this occurs, then one can take the limit as the cutoff parameters go to infinity while
readjusting the bare parameters so as to produce the desired finite answers for certain
physical quantities. In this way, one obtains finite expressions for each term in the
perturbation series expansions for all physical quantities with only the same number
of free parameters in the theory as originally were present classically. If this occurs,
the theory is called \emph{renormalizable}, and it is widely believed that physically viable
quantum field theories must be renormalizable (or, at least, satisfy properties closely
akin to renormalizability; see Weinberg 1979). Quantum electrodynamics (i.e., the
quantum field theory of a Dirac field interacting with an electromagnetic field) is
renormalizable, and the agreement to high accuracy of the predictions of the first few
terms of its perturbation series with experiments provides the best quantitative
evidence for the belief that quantum field theory provides a correct description of
nature. The Weinberg--Salam theory of the electroweak interactions and ``quantum
chromodynamics'' (i.e., the theory of quarks interacting with gluons) also are
renormalizable quantum field theories which are believed to accurately describe nature.
However, the state of affairs with regard to giving an exact formulation even of
renormalizable field theories is unsatisfactory. Only the individual terms in the
perturbation series are defined precisely, and there is good reason to believe that the
full perturbation series does not converge.

Given a classical field theory formulated in terms of a Lagrangian or Hamiltonian,
there exist a number of procedures for formulating a quantum field theory associated
with the classical theory. However, general relativity is sufficiently different from
other classical field theories that, as we shall see in more detail below, the approaches
which have been tried for formulating quantum general relativity all have encountered
fundamental difficulties. The essential difference between general relativity
and other classical theories appears to be the dual role played by the field \(g_{ab}\) as both
the quantity which describes the dynamical aspects of gravity and the quantity which
describes the background spacetime structure. Thus, it would appear that in order to
``quantize'' the dynamical degrees of freedom of the gravitational field, one must also
give a quantum mechanical description of spacetime structure. This latter problem
has no analog for other quantum field theories which are formulated on a fixed,
background spacetime, which is treated classically.

The nature of the difficulties caused by the dual role of \(g_{ab}\) are perhaps best
illustrated by the following simple example. It is a fundamental property of quantum
field theory in Minkowski spacetime that the field operator, \(\widehat\psi\), corresponding to an
integer-spin classical field, \(\psi\), evaluated at spacelike related points must commute
with itself, i.e., for \(x\) and \(x'\) spacelike related\footnote{As mentioned above, we
must in fact treat \(\widehat\psi\) as a distribution. Strictly speaking,
equation~\eqref{eq:14.1.1} should be replaced by
\([\widehat\psi(f),\widehat\psi(g)]=0\) whenever the supports of the test functions
\(f\) and \(g\) are spacelike separated.} we have
\begin{equation}
  [\widehat\psi(x),\widehat\psi(x')]
  \equiv\widehat\psi(x)\widehat\psi(x')-\widehat\psi(x')\widehat\psi(x)=0.
  \tag{14.1.1}\label{eq:14.1.1}
\end{equation}
(This equation expresses the fact that a measurement of \(\psi\) at \(x'\) cannot influence the
value of \(\psi\) at \(x\).) Now, as mentioned in section~4.4 and in problem~6 of chapter~13,
linearized gravity is just the theory of a massless, spin-2 field in Minkowski spacetime.
Thus, we may view general relativity as the theory of a self-interacting spin-2
field. By analogy with flat spacetime quantum field theories it would be natural to
expect the metric field operator \(\widehat g_{ab}\) to satisfy the commutation relation
\begin{equation}
  [\widehat g_{ab}(x),\widehat g_{cd}(x')]=0
  \tag{14.1.2}\label{eq:14.1.2}
\end{equation}
for \(x\) and \(x'\) spacelike related. However, this equation makes no sense since we do
not know if \(x\) and \(x'\) are spacelike related until we know the metric, and
equation~\eqref{eq:14.1.2} is an operator equation which, if valid, must hold independently
of the state of the gravitational field, i.e., independently of the value of (or probability
distribution for) the metric. More generally, the entire notion of causality becomes ill
defined when the notion of a classical spacetime metric is abandoned. Thus, some of
the fundamental results which are believed to hold for all other quantum field theories
appear to be inapplicable to general relativity.

The differences between general relativity and other field theories and the
difficulties with causality that arise when the spacetime metric does not have a
definite value suggest the possibility that perhaps the principles of quantum theory
do \emph{not} apply to gravity---that classical general relativity is correct at the fundamental
level. However, this viewpoint appears to be untenable, because the spacetime
metric is coupled to matter sources. Suppose spacetime structure is described by a
classical spacetime \((M,g_{ab})\) and quantum theory applies to these matter sources.
What is the curvature of spacetime associated with a given quantum state of the
matter? If the classical Einstein equation is to hold in the limit where the matter
distribution can be described classically, the most natural candidate for a quantum
version of Einstein's equation (with gravity treated classically), is\footnote{Note that
in postulating a semiclassical equation like~\eqref{eq:14.1.3}, the superposition
principle for matter states is lost, since different matter states are associated with
different spacetimes.}
\begin{equation}
  G_{ab}=8\pi\langle\widehat T_{ab}\rangle,
  \tag{14.1.3}\label{eq:14.1.3}
\end{equation}
where \(\langle\widehat T_{ab}\rangle\) denotes the expectation value of the stress-energy operator
\(\widehat T_{ab}\) in the given quantum state. Now consider a state of matter where, with
probability \(1/2\), all the matter is located in a certain region, \(O_1\), of spacetime and,
with probability \(1/2\), the matter is located in a region \(O_2\) disjoint from \(O_1\).
According to equation~\eqref{eq:14.1.3}, the gravitational field will behave like half of the matter is
in \(O_1\) and the other half is in \(O_2\). Suppose, now, that we make a measurement of the
location of the matter. We then will find the matter to be either entirely in \(O_1\) or entirely
in \(O_2\). If equation~\eqref{eq:14.1.3} continues to hold after we have resolved the quantum
state of the matter by this measurement, then the gravitational field must change in a
discontinuous, acausal manner. Thus, the attempt to treat gravity classically leads to
serious difficulties. These difficulties apparently can be avoided only by treating the
spacetime metric in a probabilistic fashion---i.e., by quantizing the gravitational field---so
that in the initial state it has probability \(1/2\) of corresponding to the gravitational field
of matter in \(O_1\), and has probability \(1/2\) of corresponding to the gravitational field
of matter in \(O_2\).

The issue of how to formulate a quantum theory of gravitation is presently under
active investigation by many researchers. We shall confine our discussion here to a
very brief mention of some of the main approaches that have been tried. A more
detailed discussion of these approaches as well as a much more complete description
of the range of topics related to quantum gravity which presently are under investigation
can be found in Isham, Penrose, and Sciama (1975, 1981) and the final
five chapters of Hawking and Israel (1979).

The \emph{covariant perturbation method} is perhaps the most straightforward approach
to formulating a quantum theory of gravity. Here one writes the spacetime metric \(g_{ab}\)
as
\begin{equation}
  g_{ab}=\eta_{ab}+\delta g_{ab},
  \tag{14.1.4}\label{eq:14.1.4}
\end{equation}
where \(\eta_{ab}\) is a flat metric and we shall assume that \(M=\mathbb R^4\) so that
\((M,\eta_{ab})\) is Minkowski spacetime. [More generally, one could replace
\((\mathbb R^4,\eta_{ab})\) by any solution, \((M,g^{(0)}_{ab})\), of Einstein's equation.]
To first order in \(\delta g_{ab}\), the classical Einstein equation is just the equation for a
free, spin-2 field. Therefore, as already mentioned above, we may view the full Einstein
equation (with \(\delta g_{ab}\) not assumed to be ``small'') as the sum of this ``free'' piece
plus a nonlinear, self-interaction term; i.e., we may view Einstein's equation as an equation
for a self-interacting spin-2 field \(\delta g_{ab}\) in Minkowski spacetime
\((\mathbb R^4,\eta_{ab})\). The covariant method treats gravity by viewing it in this manner
as an ordinary ``Poincaré covariant'' field theory. The dynamical variable,
\(\delta g_{ab}\), has considerable gauge arbitrariness, but it is known how to obtain a
perturbation series expansion for the quantum field theory of ``non-abelian'' gauge fields
of this type (Fadde'ev and Popov 1967; DeWitt 1967a,b). Thus, in this approach there is
no obstacle, in principle, to obtaining a formal perturbation series for the type of
physical quantities usually calculated for field theories in Minkowski spacetime. Note
that the flat, background metric \(\eta_{ab}\) introduced in equation~\eqref{eq:14.1.4} is
treated in an entirely classical manner in this approach.

The main difficulty which arises in this approach is that the perturbation theory one
obtains for \(\delta g_{ab}\) is non-renormalizable. Indeed, the non-renormalizability of quantum
gravity in the covariant perturbation approach can be seen purely from dimensional
arguments. The expansion parameter in the perturbation series one obtains is the
squared Planck length,\footnote{The classical vacuum Einstein equation, \(G_{ab}=0\),
does not involve Newton's constant \(G\). However, \(G\) enters the quantum theory
through the normalization of \(\delta g_{ab}\), which accounts for the presence of
\(G\) (via \(\ell_P\)) in the quantum theory.} \(\ell_P^2\). Thus, the terms in the
perturbation series of successively higher order in \(\ell_P^{2}\) must have the dimensions
of correspondingly higher powers of inverse length and the ``superficial degree of
divergence'' (see, e.g., Coleman 1973) of these terms increases. Hence, there is no
possibility of canceling the divergent terms with the ``bare'' terms. Thus, the only hope
for obtaining a well defined perturbation series without the introduction of new parameters
is that the perturbation series be \emph{finite} in each order, i.e., that when all the
contributions to order \(\ell_P^{2n}\) are added together, the divergences will cancel each
other. In fact, this happens in the lowest (``one loop'') order. However, finiteness does
\emph{not} occur in lowest order for gravity coupled to matter fields, nor is it expected to
occur in higher orders in pure quantum gravity; see Deser, van Nieuwenhuizen, and
Boulware (1975) and the references cited therein for further discussion.

Thus, in the covariant perturbation approach to formulating a quantum theory of
gravity, it appears that meaningful physical predictions cannot be made. In addition,
this approach has a number of other unappealing features. The breakup of the metric
into a background metric which is treated classically and a dynamical field
\(\delta g_{ab}\), which is quantized, is unnatural from the viewpoint of classical general
relativity. Furthermore, the perturbation theory one obtains from this approach will,
in each order, satisfy causality conditions with respect to the background metric
\(\eta_{ab}\) rather than the true metric \(g_{ab}\). Although the summed series (if it were to
converge) still could satisfy appropriate causality conditions, the covariant perturbation
approach would provide a very awkward way of displaying the role of the spacetime
metric in causal structure. Finally, in this approach it is very difficult even to formulate
questions about such issues as the quantum effects occurring near the initial singularity
of the universe, since the usual procedures for formulating quantum field theories in
Minkowski spacetime effectively assume that the interactions are ``turned off'' in the
distant past. On the other hand, the covariant approach has the advantage that one
can obtain concrete expressions for physical quantities in perturbation theory without
having to develop an entirely new conceptual framework.

An important approach which does not rely on a breakup of the spacetime metric
such as~\eqref{eq:14.1.4} is the \emph{canonical quantization method}. This approach for
formulating a quantum theory is applicable if the classical theory has been put in
Hamiltonian form. The basic idea here is (i) to take the states of the system to be
described by wave functions \(\Psi(q)\) of the configuration variables, (ii) to replace each
momentum variable by differentiation with respect to the conjugate configuration
variable, and (iii) to determine the time evolution of \(\Psi\) via the Schrödinger equation,
\(\ii\hbar\,\partial\Psi/\partial t=\widehat H\Psi\), where \(\widehat H\) is an operator corresponding
to the classical Hamiltonian \(H(p,q)\). A precise formulation of these rules for simple
quantum mechanical systems can be found in Ashtekar and Geroch (1974).

As discussed in appendix~E, general relativity can be cast in Hamiltonian form,
so one can attempt to apply the canonical quantization rules to general relativity.
However, a serious difficulty arises because of the presence of the constraint
(E.2.33). Attempts to solve this constraint (so that one can obtain configuration
variables representing only the ``true dynamical degrees of freedom'') or to impose
this constraint as an additional condition on the state vector have not been successful.
Thus, the difficulties caused by the constraint (E.2.33) in the Hamiltonian formulation
of general relativity have proven to be a serious obstacle to obtaining a
quantum theory of gravity via the canonical approach. We refer the reader to
Ashtekar and Geroch (1974) and Kuchař (1981) for further discussion.

Another important approach to formulating a quantum theory of gravity is the
\emph{path integral method}. The viewpoint here is to stress the amplitudes for physical
processes (rather than states or operators) as the fundamental entities of the theory.
For the quantum theory of a field \(\psi\) on Minkowski spacetime derived from an
action \(S[\psi]\) (see appendix~E), one expresses the amplitude for the field to go from
the function \(\psi_1(\boldsymbol{x})\) on the hypersurface \(t=t_1\) to the function
\(\psi_2(\boldsymbol{x})\) on the hypersurface \(t=t_2\) as a path integral:
\begin{equation}
  \langle\psi_1,t_1\mid\psi_2,t_2\rangle
  =\int e^{\ii S[\psi]}\,\dd\mu[\psi].
  \tag{14.1.5}\label{eq:14.1.5}
\end{equation}
Here, the integral is to be taken over \emph{all} field configurations \(\psi\) (not just those
satisfying the classical field equation) in the spacetime region between \(t_1\) and \(t_2\)
which ``interpolate'' between \(\psi_1\) and \(\psi_2\), and \(\dd\mu[\psi]\) denotes a measure
on the space of field configurations. The major difficulty which arises in the path
integral approach is the definition of the measure \(\dd\mu[\psi]\). One simply does not know
how to make mathematical sense of \(\dd\mu[\psi]\), except in the context of perturbation
theory about a free field. (In that context, one can define the integral~\eqref{eq:14.1.5}
so as to obtain the standard perturbation series, and, indeed, the path integral approach
provides a simple formal derivation of this series which is particularly useful in the case
of gauge theories; see Abers and Lee 1973.) However, despite the inability to define
\eqref{eq:14.1.5} rigorously, the formal manipulations suggested by the path integral
viewpoint have provided valuable insights\footnote{In particular, the path integral
viewpoint provides a simple explanation of how the classical limit arises. Namely,
in appropriate circumstances the dominant contributions to the integral~\eqref{eq:14.1.5}
should arise from field configurations where the phase factor \(S[\psi]\) is extremized,
since then the nearby configurations will add coherently. But these field configurations
are precisely the solutions of the classical equations (see appendix~E).} and useful
approximation schemes.

General relativity can be derived from an action principle (see appendix~E), so one
can attempt to formulate a quantum theory of gravity by the path integral approach.
It would be natural to write down an integral of the form
\begin{equation}
  Z=\int e^{\ii S[g_{ab}]}\,\dd\mu[g_{ab}]
  \tag{14.1.6}\label{eq:14.1.6}
\end{equation}
with \(S\) given by equation (E.1.13) or equation (E.1.42), to represent the amplitude,
\(\langle(h_1)_{ab},t_1\mid(h_2)_{ab},t_2\rangle\), for going from spatial metric
\((h_1)_{ab}\) at time \(t_1\) to \((h_2)_{ab}\) at time \(t_2\). However, in general
relativity the ``times'' \(t_1\) and \(t_2\) are merely coordinate labels and do not have
physical significance. Thus, the amplitude
\(\langle(h_1)_{ab},t_1\mid(h_2)_{ab},t_2\rangle\) is not a physically meaningful
quantity. The underlying reason for this difficulty is closely related to the difficulty
of the canonical quantization approach: The ``true dynamical degrees of freedom'' of
the gravitational field have not been isolated. As mentioned near the end of appendix~E,
general relativity is much like a ``parameterized theory'' where time is treated as a
dynamical variable. Thus, it is redundant to insert the ``label times'' \(t_1\) and \(t_2\)
in the amplitudes, as, in effect, an ``intrinsic time'' already is present in the canonical
variables \(h_{ab}\), \(\pi^{ab}\) describing the gravitational field. However, since one
does not have a well defined decomposition of these variables into ``true dynamical
degrees of freedom'' and ``intrinsic time variables,'' it is far from clear precisely what
physical amplitude \(Z\) is supposed to represent. Furthermore, another fundamental
difficulty which arises is the usual problem in the path integral approach of giving
precise meaning to \(\dd\mu[g_{ab}]\) so that the right-hand side of
equation~\eqref{eq:14.1.6} is well defined. On the other hand, in the path integral
approach one can envision analyzing issues such as the probability for a change of
spatial topology by including in the integral~\eqref{eq:14.1.6} spacetime metrics for
which such a spatial topology change occurs. It is difficult to see how this issue could
even be formulated in the canonical approach (since the presence of a Hamiltonian
requires, in essence, global hyperbolicity, and hence no change of spatial topology) or
in the covariant approach.

An important variant of the above path integral approach is the \emph{Euclidean path
integral approach}. In field theories in Minkowski spacetime, a number of quantities
which arise---in particular the vacuum expectation values of products of field operators---
are holomorphic functions of the global inertial coordinates \(t,x,y,z\) in a
domain that includes negative imaginary values of the time coordinate, i.e.,
\(t=-\ii\tau\), where \(\tau\) is real and positive (see Streater and Wightman 1964). It is
useful to view this type of analytic continuation in the following manner. We define
\emph{complexified Minkowski spacetime} to be the four-complex-dimensional manifold
\(\mathbb C^4\) (which also may be viewed as an eight-real-dimensional manifold) with
complex metric \(\eta_{ab}\) defined in terms of the complex Cartesian coordinates
\(t,x,y,z\) of \(\mathbb C^4\) by
\begin{equation}
  \dd s^2=-\dd t^2+\dd x^2+\dd y^2+\dd z^2.
  \tag{14.1.7}\label{eq:14.1.7}
\end{equation}
Thus by restricting to real values of \(t,x,y,z\), we recover ordinary Minkowski
spacetime as a four-real-dimensional submanifold of \((\mathbb C^4,\eta_{ab})\). However,
by restricting \(x,y,z\) to be real but \(t\) to be pure imaginary, we obtain another
four-real-dimensional submanifold, now with a real, Euclidean metric
\begin{equation}
  \dd s^2=+\dd\tau^2+\dd x^2+\dd y^2+\dd z^2
  \tag{14.1.8}\label{eq:14.1.8}
\end{equation}
where \(\tau=\ii t\). This submanifold is referred to as a \emph{Euclidean section} of
\((\mathbb C^4,\eta_{ab})\). Thus, we may view the above analytic continuation of
functions to negative imaginary values of \(t\) as the evaluation of these functions on
the positive \(\tau\) region of a Euclidean section of complexified Minkowski spacetime.
We may perform our analysis of the field theory in this Euclidean section and then
analytically continue the functions back to the Minkowskian section to obtain the
physical predictions. In the path integral approach, there is considerable potential
advantage to proceeding in this manner, since in many theories the Euclidean action,
\(S_E[\psi]=-\ii S[\psi]\rvert_{t=-\ii\tau}\), is positive definite and the analytically
continued integral is exponentially damped at ``large \(\psi\),'' rather than oscillatory.
Thus, the possibility of making mathematical sense of the integral in
equation~\eqref{eq:14.1.5} appears to be greatly enhanced in the Euclidean section.

Many difficulties arise when one attempts to apply the Euclidean path integral
approach to quantum gravity. In general relativity, one does not have a natural
background flat spacetime \((\mathbb R^4,\eta_{ab})\) which can be ``complexified'' to
allow analytic continuation to be performed. Nevertheless, one could try to analytically
continue Lorentzian signature metrics to Riemannian metrics and work with a path
integral over Riemannian metrics. However, except in special cases such as static
spacetimes, it is generally impossible to represent an analytic spacetime \((M,g_{ab})\)
as a ``Lorentzian section'' of a four-complex-dimensional manifold with complex metric
which possesses a ``Euclidean section,'' i.e., a four-real-dimensional submanifold with
real, Riemannian metric. Thus, one does not have a general prescription for analytically
continuing Lorentz signature metrics to Riemannian metrics. (Furthermore, even if
one did, one does not have any theorems guaranteeing the analyticity of any quantities
arising in quantum gravity.) Nevertheless, one can postpone the resolution of
interpretational issues and study the properties of the path integral~\eqref{eq:14.1.6}
over Riemannian metrics, with \(iS[g_{ab}]\) replaced by \(-S_E[g_{ab}]\), where \(S_E\)
is given by the analog of equation (E.1.42) for Riemannian metrics.\footnote{The action
\(S_E\) is not positive definite. However, it is positive definite for asymptotically
Euclidean metrics with \(R=0\) (Schoen and Yau 1979), and Hawking (1979) has
proposed a further analytic continuation of the ``conformal degree of freedom'' of the
metric which, in effect, makes \(S_E\) positive definite.} A number of highly suggestive
results have been obtained in this manner (Hawking 1979). Perhaps the most dramatic
success of the Euclidean approach is that, as explained in section~14.3, the creation of
particles by a Schwarzschild black hole can be related in a direct and simple manner
to properties of the Euclidean Schwarzschild solution.

A considerably more unconventional approach toward the formulation of a quantum
theory of gravity is provided by the \emph{twistor approach}. A \emph{twistor} in Minkowski
spacetime may be defined as a pair,
\begin{equation}
  Z=(\omega^A,\pi_{A'}),
  \tag{14.1.9}\label{eq:14.1.9}
\end{equation}
consisting of a spinor field, \(\omega^A\), and a complex conjugate spinor field
\(\pi_{A'}\), satisfying the \emph{twistor equation},
\begin{equation}
  \partial_{AA'}\omega^B=-\ii\epsilon_A{}^B\pi_{A'}.
  \tag{14.1.10}\label{eq:14.1.10}
\end{equation}
We refer the reader to such references as Penrose (1967), Penrose and MacCallum
(1972), Penrose (1975), and Penrose and Ward (1980) for the motivation for
introducing twistors and a discussion of their properties. The collection of twistors on
Minkowski spacetime forms a four-complex-dimensional vector space, and the
projective twistors---i.e., the equivalence classes of twistors which differ by a nonzero
complex multiple---comprise the three-complex-dimensional manifold \(\mathbb{CP}^3\).
There exist a number of natural correspondences between Minkowski spacetime and
twistor space. For example, a null geodesic in Minkowski spacetime corresponds to
a point in null projective twistor space (i.e., the space of projective twistors satisfying
\(Z\mathbin{\cdot}\overline Z\equiv
\omega^A\overline\pi_A+\overline\omega^{A'}\pi_{A'}=0\)), whereas a point in
Minkowski spacetime corresponds to a 2-sphere of null projective twistors. Further
correspondences between compactified, complexified Minkowski spacetime and twistor
space also may be obtained. Furthermore, certain differential equations in Minkowski
spacetime can be reformulated as analyticity conditions in twistor space (see Ward 1981).

The basic starting point of the twistor quantization approach is to use twistor space
rather than spacetime as the underlying classical manifold structure upon which the
quantum fields are defined. In view of the correspondences between spacetime and
twistor space, this should result in a quantum theory where, roughly speaking,
certain causal relationships in spacetime retain their classical properties, but the
notion of spacetime points becomes ``fuzzy.'' (In essentially all other approaches, the
notion of spacetime points is well defined but, as discussed above, since the spacetime
metric is described probabilistically, the notion of causality becomes ``fuzzy.'')
Hence, in the twistor approach one can envision incorporating causality into the
quantum theory in a natural way. Thus far, twistors have proven to be useful
mathematical tools (see, e.g., Atiyah and Ward 1977). However, at present, little
progress has been made toward formulating a concrete quantum theory of gravity via
the twistor approach.

Given the present lack of success of the above approaches to formulating a
quantum version of general relativity, it is natural to consider modifications of
classical general relativity which may lead to better behavior of the quantum theory.
In particular, one may seek to modify the Einstein field equation so that the quantum
theory becomes renormalizable in the covariant perturbation approach. Perhaps the
simplest attempt along these lines is to modify the Einstein Lagrangian (E.1.12) by
adding terms quadratic in the curvature. The coordinate component form of the new
field equation then involves fourth derivatives of the metric, so this theory often is
referred to as a ``higher derivative'' theory of gravity. Stelle (1977) has shown that
the quantum version of this theory is formally renormalizable. However, other
serious difficulties arise in the theory, so it does not appear that this theory is
physically viable.

A much more ambitious attempt to modify general relativity is given by \emph{supergravity
theories}. Supergravity theories are a class of models involving interacting
fields of spin 0, \(1/2\), 1, \(3/2\), and 2. (The prefix ``super'' refers to the fact that these
models have a certain type of symmetry called ``supersymmetry'' between the boson
[i.e., integer spin] and fermion [i.e., half-integer spin] degrees of freedom.) The
main goals of supergravity theories are (i) to improve the renormalizability (or
finiteness) properties of quantum gravity and (ii) to unify gravity with the other basic
interactions of nature. With regard to the first goal, it has been shown that supergravity
is finite at least to ``two loop'' order in perturbation theory. At present, it is
not known how far in perturbation theory this finiteness extends, and it is believed
to be possible that supergravity could be finite to all orders. The second goal is
particularly ambitious since one asks supergravity to account for all interactions
observed in high energy particle physics. At present, there appear to be serious
difficulties in reaching this goal. We refer the reader to van Nieuwenhuizen (1981)
and the articles in Hawking and Roček (1981) for an introduction to supergravity,
a discussion of some recent research, and references to earlier work on supergravity
and supersymmetry.

In summary, although many approaches have been tried, there presently does not
exist a demonstrably viable quantum theory of gravity. It is possible that the
difficulties are basically technical in nature and that, for example, a better procedure
for dealing with the constraint (E.2.33) will lead to a satisfactory formulation of
quantum gravity via the canonical approach, or that supergravity will be shown to
be a fully satisfactory theory. Alternatively, it is possible that the difficulties may be
of a very fundamental nature, and that the quantum theory of gravity simply does not
fit into the framework established for other quantum theories.

However, the lack of a satisfactory quantum theory of gravity does not mean that
we cannot perform any reliable calculations of quantum effects occurring in strong
gravitational fields. In atomic physics, in appropriate circumstances one can reliably
calculate electromagnetically induced transition rates of electrons in atoms using a
classical treatment of the electromagnetic field. Similarly, in quantum field theory,
using a classical treatment of the electromagnetic field one can, in appropriate
circumstances, reliably calculate the spontaneous creation of electron--positron pairs
in a strong electric field. Thus, it is expected that by treating gravity in the classical
framework of general relativity, it still should be possible to reliably calculate some
of the quantum effects that gravity produces on other fields, such as the creation of
particle--antiparticle pairs. It is to the study of these effects that we now turn.
